\documentclass[11pt,a4paper]{article}
\usepackage[utf8]{inputenc}
\usepackage[T1]{fontenc}

% Set margins without geometry package to avoid miktex issues
\setlength{\textwidth}{6.5in}
\setlength{\textheight}{9.0in}
\setlength{\topmargin}{-0.5in}
\setlength{\oddsidemargin}{0in}
\setlength{\evensidemargin}{0in}
\setlength{\parindent}{0pt}
\setlength{\parskip}{0.5em}

\title{\textbf{BrainBlock DRL Project}\\ \large Design Choices \& Task Plan}
\author{Group Plan}
\date{\today}

\begin{document}
\maketitle

\section{Design Decisions (with Pros \& Cons)}

\subsection{Invalid Action Handling}
When the agent selects an action that results in an illegal placement (out-of-bounds or overlap):

\vspace{0.5em}
\noindent
\begin{tabular}{|l|p{3.5cm}|p{3.5cm}|p{3.5cm}|}
\hline
\textbf{Option} & \textbf{Description} & \textbf{Pros} & \textbf{Cons} \\
\hline
\textbf{A) Hard term} & Episode ends immediately on any invalid action & Simple to implement; learns to avoid invalid quickly & Few steps early in training $\rightarrow$ slow learning \\
\hline
\textbf{B) Neg reward + skip} & Give negative reward, skip turn, keep piece & Keeps playing $\rightarrow$ more learning signal; forgiving & Agent may ``farm'' negative rewards; longer episodes \\
\hline
\textbf{C) Action masking} & Mask out invalid actions & No wasted steps; fastest learning & Complex implementation; must compute mask \\
\hline
\textbf{D) Neg reward + retry} & Agent retries until valid & Good exploration & Can get stuck in loops; adds hyperparameter \\
\hline
\end{tabular}

\textbf{Recommendation:} Start with \textbf{C (action masking)} for the main experiments. It is the cleanest and most efficient. Optionally compare with \textbf{A} as a second reward-function experiment.

\subsection{State Representation}

\noindent
\begin{tabular}{|l|p{3.5cm}|p{3.5cm}|p{3.5cm}|}
\hline
\textbf{Option} & \textbf{Description} & \textbf{Pros} & \textbf{Cons} \\
\hline
\textbf{A) Multi-channel binary} & Ch0: board. Ch1--5: one-hot piece. Ch6+: remaining & Spatial structure preserved; works well with CNNs & Small board (8x5) means CNN may be overkill \\
\hline
\textbf{B) Flat vector} & Flatten everything to 50-dim vector & Simple; works well with MLPs; easy to debug & Loses spatial relationships between cells \\
\hline
\textbf{C) Grid + inventory} & Board as channel, shape rendered, counts broadcast & Richest spatial info; strong CNN support & Largest observation size; slightly more complex \\
\hline
\end{tabular}

\textbf{Recommendation:} Both \textbf{A} and \textbf{C} are excellent choices. As you noted (\textit{``belki olabilir''}), \textbf{C (Multi-channel grid + inventory channel)} is mathematically richer and a very strong choice if you want to maximize the CNN's spatial awareness. Option \textbf{A} is a solid standard starting point.

\subsection{Piece Orientation Encoding}

\noindent
\begin{tabular}{|l|p{3.5cm}|p{3.5cm}|p{3.5cm}|}
\hline
\textbf{Option} & \textbf{Description} & \textbf{Pros} & \textbf{Cons} \\
\hline
\textbf{A) All 8 orientations} & Every piece has actions for all 8 orientations & Uniform action space (320 actions); simple & Redundant actions waste exploration \\
\hline
\textbf{B) Map redundant} & Internally deduplicate combinations & No wasted action slots; cleaner & Interface still has 320 actions \\
\hline
\end{tabular}

\textbf{Recommendation:} Use \textbf{A} for the action interface but implement \textbf{B} internally --- mask out redundant orientation indices so the agent never picks them.

\subsection{Reward Function Designs}

\textbf{Reward Function 1: Sparse Completion Reward}
\begin{itemize}
    \item Successful piece placement: $+1$
    \item Board fully solved (all 10 pieces): $+100$ bonus
    \item Invalid action: $-10$, episode ends
    \item Dead-end: $0$, episode ends
\end{itemize}
\textit{Pros:} Simple. \textit{Cons:} Sparse; slow convergence.

\textbf{Reward Function 2: Dense Shaped Reward}
\begin{itemize}
    \item Successful piece placement: $+4$ (= cells covered)
    \item Complete row: $+5$ bonus
    \item Contiguous empty region: $+2$ bonus
    \item Board fully solved: $+50$ bonus
    \item Invalid action: $-5$, episode ends
    \item Dead-end: $-$ (remaining empty cells)
\end{itemize}
\textit{Pros:} Rich gradient signal. \textit{Cons:} May bias the agent toward suboptimal strategies.

\textbf{Recommendation:} Compare \textbf{Reward 1 (sparse)} vs. \textbf{Reward 2 (dense shaped)}.

\subsection{DRL Algorithm Choice}

\noindent
\begin{tabular}{|l|p{3.5cm}|p{3.5cm}|p{3.5cm}|}
\hline
\textbf{Option} & \textbf{Description} & \textbf{Pros} & \textbf{Cons} \\
\hline
\textbf{A) DQN} & Value-based; discrete action fit & Well-understood; simple & Slow convergence; no stochastic policy \\
\hline
\textbf{B) PPO} & Policy gradient with clipped surrogate & State-of-the-art; stable; naturally stochastic & More complex to implement \\
\hline
\textbf{C) A2C} & Synchronous actor-critic & Simpler than PPO & Less stable; higher variance \\
\hline
\end{tabular}

\textbf{Recommendation:} Use \textbf{PPO} as the primary algorithm because its stochastic policy naturally discovers multiple solutions easily.

\subsection{Network Architecture}

\noindent
\begin{tabular}{|l|p{3.5cm}|p{3.5cm}|p{3.5cm}|}
\hline
\textbf{Option} & \textbf{Description} & \textbf{Pros} & \textbf{Cons} \\
\hline
\textbf{A) MLP only} & Flatten observation $\rightarrow$ 2-3 hidden layers & Simple; fast & Ignores spatial structure \\
\hline
\textbf{B) CNN + MLP head} & 2-3 conv layers $\rightarrow$ flatten $\rightarrow$ MLP & Captures spatial patterns; good for grid & Board is very small \\
\hline
\textbf{C) CNN + params} & CNN on board, MLP on inventory vector & Best info separation & Most complex; overkill \\
\hline
\end{tabular}

\textbf{Recommendation:} Start with \textbf{B (CNN + MLP head)} --- it's a good default for grid-based environments.

\subsection{Episode Termination Conditions}

\noindent
\begin{tabular}{|l|p{3.5cm}|p{7.5cm}|}
\hline
\textbf{Condition} & \textbf{When} & \textbf{Notes} \\
\hline
\textbf{Success} & All 10 pieces placed, board covered & Return large positive reward \\
\hline
\textbf{Invalid action} & Agent picks out-of-bounds or overlapping & Only relevant if NOT using action masking \\
\hline
\textbf{Dead-end} & No valid actions exist for current piece & Episode must end; shaped reward penalty \\
\hline
\textbf{Max steps} & Count exceeds a limit (e.g., 50) & Safety net; unlikely if using masking \\
\hline
\end{tabular}

\subsection{Piece Queue: Visible vs. Hidden}

\noindent
\begin{tabular}{|l|p{3.5cm}|p{3.5cm}|p{3.5cm}|}
\hline
\textbf{Option} & \textbf{Description} & \textbf{Pros} & \textbf{Cons} \\
\hline
\textbf{A) Full queue visible} & Agent sees current piece + all remaining order & Agent can plan ahead; richer state & Much larger state space \\
\hline
\textbf{B) Hidden queue} & Agent sees current piece + remaining counts only & Simpler state; matches spec & Cannot plan for specific orderings \\
\hline
\end{tabular}

\textbf{Recommendation:} Use \textbf{B} --- the spec requires "information about remaining pieces" but doesn't require order. Counts are simpler and more generalizable.

\subsection{Coordinate Convention}

\noindent
\begin{tabular}{|l|p{11.0cm}|}
\hline
\textbf{Option} & \textbf{Description} \\
\hline
\textbf{A) Row-major} & y-axis = row (top/bot), x-axis = col (L/R). Action = (orient, col, row). \\
\hline
\textbf{B) Cartesian} & x = col (L/R), y = row (bot/top). Action = (orient, x, y). \\
\hline
\end{tabular}

\textbf{Recommendation:} Use \textbf{A (row-major)} --- it's the NumPy convention.

\section{Detailed Task List — Split for 2 Members}

\subsection{Ground Rules}
\begin{itemize}
    \item Each member builds their \textbf{own complete pipeline} (environment + agent + training + eval).
    \item This ensures both can work fully in parallel without blocking each other.
    \item Common interfaces and hyperparameters are agreed upon upfront.
    \item At the end, combine the best results and write the report together.
\end{itemize}

\subsection{Shared Agreements (Discuss Before Starting)}
\begin{itemize}
    \item[ [ ] ] Agree on state representation (Option A or C).
    \item[ [ ] ] Agree on action space: 320 flat discrete actions (8 orientations x 8 x 5).
    \item[ [ ] ] Agree on coordinate convention: row-major.
    \item[ [ ] ] Agree on piece shapes (offsets for each orientation of I, O, L, Z, T).
    \item[ [ ] ] Agree on random seeds to use: [42, 123, 456, 789, 1024].
    \item[ [ ] ] Set up shared Git repo with folder structure.
\end{itemize}

\subsection{Suggested Folder Structure}
\begin{verbatim}
cs545-project/
|-- common/                  # Shared utilities (pieces, visualize)
|   |-- pieces.py            # Piece shapes and orientations
|   |-- visualize.py         # Board visualization utility
|-- member_A/                # Member A's full pipeline (sparse reward)
|   |-- environment.py       
|   |-- agent.py             
|   |-- network.py           
|   |-- train.py             
|   |-- evaluate.py          
|   |-- configs/             
|-- member_B/                # Member B's full pipeline (dense reward)
|   |-- environment.py       
|   |-- agent.py             
|   |-- network.py           
|   |-- train.py             
|   |-- evaluate.py          
|   |-- configs/             
|-- results/                 # Combined results for report
|-- report/                  # LaTeX report
|-- presentation/            # Slides
\end{verbatim}

\subsection{Member A — Tasks}
\textbf{Phase 1: Core Environment (Days 1--3)}
\begin{itemize}
    \item[ [ ] ] Implement piece definitions (I, O, L, Z, T).
    \item[ [ ] ] Implement BrainBlockEnv with \textbf{Sparse Reward}.
    \item[ [ ] ] Write unit tests for environment.
    \item[ [ ] ] Implement board visualization.
\end{itemize}
\textbf{Phase 2: Agent \& Training (Days 4--7)}
\begin{itemize}
    \item[ [ ] ] Implement PPO agent from scratch in PyTorch.
    \item[ [ ] ] Implement training loop (rollouts, updates, logging).
    \item[ [ ] ] Implement configurable hyperparameters.
\end{itemize}
\textbf{Phase 3: Experiments \& Evaluation (Days 8--10)}
\begin{itemize}
    \item[ [ ] ] Train with 5 random seeds.
    \item[ [ ] ] Implement evaluation script (success rate, return, length).
    \item[ [ ] ] Generate learning curve plots.
    \item[ [ ] ] Collect 5 distinct solutions.
\end{itemize}
\textbf{Phase 4: Report Sections (Days 11--12)}
\begin{itemize}
    \item[ [ ] ] Write: MDP formulation, sparse reward explanation, and algorithms.
\end{itemize}

\subsection{Member B — Tasks}
\textbf{Phase 1: Core Environment (Days 1--3)}
\begin{itemize}
    \item[ [ ] ] Implement piece definitions (I, O, L, Z, T).
    \item[ [ ] ] Implement BrainBlockEnv with \textbf{Dense Shaped Reward}.
    \item[ [ ] ] Write unit tests for environment.
\end{itemize}
\textbf{Phase 2: Agent \& Training (Days 4--7)}
\begin{itemize}
    \item[ [ ] ] Implement PPO agent (or DQN) for fair comparison.
    \item[ [ ] ] Implement training loop.
\end{itemize}
\textbf{Phase 3: Experiments \& Evaluation (Days 8--10)}
\begin{itemize}
    \item[ [ ] ] Train with 5 random seeds.
    \item[ [ ] ] Generate learning curve plots.
    \item[ [ ] ] Collect 5 distinct solutions + qualitative analysis.
\end{itemize}
\textbf{Phase 4: Report Sections (Days 11--12)}
\begin{itemize}
    \item[ [ ] ] Write: Dense reward logic, comparative failure analysis, conclusions.
\end{itemize}

\subsection{Joint Tasks (Days 12--14)}
\begin{itemize}
    \item[ [ ] ] Merge results and create side-by-side comparative figures.
    \item[ [ ] ] Finalize the report PDF.
    \item[ [ ] ] Prepare the 10-minute presentation slides.
    \item[ [ ] ] Prepare live visual demo script.
\end{itemize}

\section{Key Hyperparameters to Start With}
\begin{tabular}{|l|l|}
\hline
\textbf{Parameter} & \textbf{Suggested Value} \\
\hline
Learning rate & 3e-4 \\
Discount ($\gamma$) & 0.99 \\
GAE $\lambda$ & 0.95 \\
PPO clip $\epsilon$ & 0.2 \\
Entropy coefficient & 0.01 \\
Value loss coefficient & 0.5 \\
Max grad norm & 0.5 \\
Rollout steps & 2048 \\
Mini-batch size & 64 \\
PPO epochs per update & 4 \\
Total training episodes & 500,000+ \\
Network hidden dim & 256 \\
CNN channels & [32, 64] \\
CNN kernel size & 3x3 with padding \\
\hline
\end{tabular}

\end{document}
